% reference-mark-interlinear.tex -- V3 fixture
% reference_mark=interlinear under the pdfTeX backend. Asserts that an
% interlinear footnote mark does not disturb the spacing that the pTeX
% route computes via \lastnodechar + \jlreq@fixjfm (penalty/JFM glue
% between the char before the mark and the char after it):
%   A) \hbox{あ、<mark>あ} must have the same natural width as
%      \hbox{あ、あ} (pair glue 、->阿 from the class table; the
%      zero-width mark itself adds nothing)
%   B) \hbox{あ<mark>a} must have the same natural width as \hbox{あa}
%      (CJK->LA boundary via postnote -> \ptj@peek@a; pTeX would
%      insert the same xkanjiskip through \jlreq@fixjfm)
%   C) a paragraph full of interlinear marks must still typeset (rc=0)
\documentclass[pdflatex,paper=b5]{jlreq}
\jlreqsetup{reference_mark=interlinear}
\begin{document}
\makeatletter
\newbox\inter@boxa
\newbox\inter@boxb
% 計測間の汚染（前ボックス末尾の欧文が LA->JA 境界 glue を次ボックス先頭に
% 作る）を防ぐため，各計測は直前状態を中立化してから行う．
\def\inter@resetstate{\global\ptj@lastclass=90\relax\global\ptj@pend=0\relax}
\inter@resetstate
\setbox\inter@boxa=\hbox{あ、\footnote{注1}あ}
\inter@resetstate
\setbox\inter@boxb=\hbox{あ、あ}
\@tempdima=\dimexpr\wd\inter@boxa - \wd\inter@boxb\relax
\ifdim\@tempdima<\z@ \@tempdima=-\@tempdima\fi
\ifdim\@tempdima<0.5pt
  \typeout{INTER-PASS: JA pair glue across mark = direct (\the\wd\inter@boxa\space vs \the\wd\inter@boxb)}%
\else
  \typeout{INTER-FAIL: JA pair glue across mark differs (\the\wd\inter@boxa\space vs \the\wd\inter@boxb, diff \the\@tempdima)}%
\fi
\inter@resetstate
\setbox\inter@boxa=\hbox{あ\footnote{注2}abc}
\inter@resetstate
\setbox\inter@boxb=\hbox{あabc}
\@tempdima=\dimexpr\wd\inter@boxa - \wd\inter@boxb\relax
\ifdim\@tempdima<\z@ \@tempdima=-\@tempdima\fi
\ifdim\@tempdima<0.5pt
  \typeout{INTER-PASS: CJK->LA boundary across mark = direct (\the\wd\inter@boxa\space vs \the\wd\inter@boxb)}%
\else
  \typeout{INTER-FAIL: CJK->LA boundary across mark differs (\the\wd\inter@boxa\space vs \the\wd\inter@boxb, diff \the\@tempdima)}%
\fi
% B2) 状態連続性の直接検証: 合印（と脚注本文処理）の中で汚れる
%    \ptj@lastclass/\ptj@pend が postnote で正しく復帰し，
%    「そこから見える直前クラス」が脚注手前の和文クラス（漢=0）に戻ること．
\setbox\inter@boxa=\hbox{漢\footnote{注3}%
  \ifnum\ptj@lastclass=0\relax
    \ifnum\ptj@pend=0\relax
      \typeout{INTER-PASS: class/pend restored across mark (lc=\the\ptj@lastclass, pend=\the\ptj@pend)}%
    \else
      \typeout{INTER-FAIL: pend=\the\ptj@pend\space after restore (lc=\the\ptj@lastclass)}%
    \fi
  \else
    \typeout{INTER-FAIL: lc=\the\ptj@lastclass\space after restore (pend=\the\ptj@pend)}%
  \fi
  abc}
\noindent
以下は合印の連続使用（組版成立の確認）。テスト用の文です\footnote{あいうえお}、
句読点の直後にも合印が来ます\footnote{かきくけこ}。欧文への接続abc\footnote{数字123}def、
数式への接続\(x+y\)\footnote{さしすせ}たちつてと。漢字漢字漢字漢字漢字漢字漢字漢字漢字漢字
漢字漢字漢字漢字漢字漢字漢字漢字\footnote{折返}漢字漢字漢字漢字漢字漢字漢字漢字漢字漢字。
\end{document}
